\section{从人工写入转向运行中的字节接收}

上一章的操作人员可以在复位期间逐地址写 RAM。机器一旦离开实验台，这种方式就出现了
具体成本：每次更换程序都要停机、连接调试工具、重新写入多个区域，并由人保证最后一个
字节已经完成。机器自身无法重复这套过程。

一种选择是把永远不变的程序直接固化在启动存储中。这样不再需要每次装入，但更换程序要
重新写非易失器件；输入参数和可写栈仍要另行准备。本章需要的是“保持固定启动程序不变，
运行期间接收另一份可替换程序”，所以必须给运行中的处理器增加一个能够取得外部字节的
接口。

\subsection{为什么不能收到一条指令就立刻执行}

外部连接按时间送来字节，处理器执行程序却会分支回到先前地址。若机器收到四个指令字节就
立即执行，循环第二次取指时，先前字节已经从顺序数据流中消失。除非外部端同时承担一片可
随机访问的指令存储，否则流式执行无法支持本书的求和循环。

另一个问题是完整性。前几条指令到达后便开始执行，后续传输若中断或校验失败，已经执行的
指令可能改写 RAM。此时不能再把整次接收当作“从未发生”。因此，基本方案先把全部程序
写入 RAM，验证通过以后才改变处理器入口。

\begin{longtable}{|L{0.24\linewidth}|L{0.28\linewidth}|L{0.34\linewidth}|}
\hline
\rowcolor{BookTableHead}
方案 & 省去的步骤 & 失去的能力或边界 \\ \hline
程序固化在启动存储 & 每次传输与装入 & 不能方便替换程序 \\ \hline
收到一条执行一条 & 完整 RAM 映像 & 不能可靠分支回旧指令 \\ \hline
接收与执行重叠 & 等待完整传输的时间 & 不能在产生效果前完成整体验证 \\ \hline
先接收、验证、再移交 & 无 & 需要接收状态、RAM 区域和明确提交点 \\ \hline
\end{longtable}

最后一行没有宣称适用于所有机器。它只是给本章提供一条可检查的基本路径：外部字节先成为
完整映像，处理器随后一次性从启动程序切到新入口。串行控制器只负责把线上位流整理成字节；
映像格式、范围检查和最终移交仍由启动程序负责。
